\documentclass[10pt,a4paper]{amsart}
\usepackage[margin=19mm]{geometry}
\usepackage{amsmath,amssymb,mathtools}
\usepackage[scr=rsfs,cal=euler]{mathalfa}
\usepackage{xfrac}
\usepackage[unicode,hidelinks,bookmarks=false]{hyperref}
\newcommand{\ie}{i.e.\ }
\newcommand{\cf}{cf.\ }
\newcommand{\resp}{respectively\ }
\newcommand{\Subsection}{\subsection}
\providecommand{\nobreakdash}{\nobreak\hbox{-}}
\setlength{\emergencystretch}{2em}
\multlinegap0pt
\usepackage{amssymb}
\usepackage{mathtools}
\usepackage[shortlabels]{enumitem}
\usepackage{xcolor}
\newtheorem{proposition}{Proposition}
\newcommand{\T}{\mathbb{T}}
\title{\textbf{Towards a Gagliardo-Type Theory of Fractional Sobolev Spaces on Arbitrary Time Scales}}
\author{Hafida Abbas, Abdelhalim Azzouz, Praveen Agarwal, Delfim F. M. Torres}
\date{Source: arXiv:2603.14715v2 (CC BY 4.0)}
\begin{document}
\maketitle

\noindent\emph{Source and licence.} \textbf{Towards a Gagliardo-Type Theory of Fractional Sobolev Spaces on Arbitrary Time Scales}. Version arXiv:2603.14715v2 (CC BY 4.0), distributed by arXiv under the Creative Commons Attribution 4.0 International licence, http://creativecommons.org/licenses/by/4.0/, as recorded in the arXiv licence metadata for this exact version and shown on the abstract page https://arxiv.org/abs/2603.14715v2. Full licence notice: \texttt{licenses/arxiv-2603.14715-CC-BY-4.0.txt}.\par\smallskip

\noindent\textit{Editorial scope.} Counted unit: the article's proposition prop:two-sided-poincare (source0.tex lines 1613--1626). The article's complete original proof of it is delivered with it (source0.tex lines 1628--1703, one \texttt{proof} environment of 76 lines). No mathematics is written, repaired or re-derived: the statement and the proof are the article's own lines, in the article's own notation, and no retained line is substituted or re-flowed.  Retained uncounted as the source-local context the proof needs: lines 1514--1538; lines 1531--1533.  Nothing was omitted from any retained span, so the package carries no omission record.  No retained line contains a citation, so the wrapper carries no bibliography and no bibliography entry.  No retained line contains a figure environment, an image-inclusion command or drawing code, so no image is delivered and the package declares no asset dependency.  Editorial formatting only, in addition: an AMS \texttt{amsart} preamble that restates the article's own declarations and macros used by the retained lines, namely the environments \texttt{proposition} on separate counters, together with the standard packages those lines need.  The retained environments print in this document as Proposition 1 in order of appearance, whereas the article prints the same items with the numbering of its own full document.  The complete native source of the article as distributed in the arXiv e-print for this exact version is bundled unchanged as \texttt{sources/196427/source0.tex}.
\begin{proposition}[Dependence of the Poincar\'e constant on the gap]
\label{prop:poincare-gap}
Let $\alpha\in(0,1)$, $p=2$, and
$\T_\delta=[0,1]\cup[1+\delta,2+\delta]$ for $\delta>0$. Define
\[
K(\delta,\alpha) :=
\begin{cases}
\dfrac{2(1+\delta)^{1-2\alpha}-\delta^{1-2\alpha}-(2+\delta)^{1-2\alpha}}
{2\alpha(1-2\alpha)}
& \text{if }\alpha\neq\tfrac{1}{2},\\[8pt]
\ln\!\left(\dfrac{(1+\delta)^2}{\delta(2+\delta)}\right)
& \text{if }\alpha=\tfrac{1}{2}.
\end{cases}
\]
Then the optimal Poincar\'e constant in
$\|u-u_{\T_\delta}\|_{L^2_\Delta}\le C_P(\delta)[u]_{W^{\alpha,2}_\Delta(\T_\delta)}$
satisfies
\begin{equation}\label{eq:poincare-lower}
C_P(\delta) \ge \frac{1}{2\,K(\delta,\alpha)^{1/2}}.
\end{equation}
Moreover, as $\delta\to\infty$,
\begin{equation}\label{eq:poincare-blowup}
C_P(\delta) \ge \frac{1}{2}\,\delta^{(1+2\alpha)/2}\bigl(1+o(1)\bigr).
\end{equation}
\end{proposition}

\begin{equation}
C_P(\delta) \ge \frac{1}{2\,K(\delta,\alpha)^{1/2}}.
\end{equation}

\begin{proposition}[Two-sided gap estimate for the Poincar\'e constant]
\label{prop:two-sided-poincare}
In the setting of Proposition~\ref{prop:poincare-gap} ($p = 2$, 
$\mathbb{T}_\delta = [0,1] \cup [1+\delta, 2+\delta]$), there exists a 
constant $C^* = C^*(\alpha) > 0$ such that, for every $\delta \ge 1$,
\begin{equation}\label{eq:two-sided-poincare}
\frac{1}{2 \, K(\delta, \alpha)^{1/2}} 
\;\le\; C_P(\delta) 
\;\le\; C^* \, \delta^{(1+2\alpha)/2}.
\end{equation}
In particular, $C_P(\delta) \asymp \delta^{(1+2\alpha)/2}$ as 
$\delta \to \infty$, and the lower bound \eqref{eq:poincare-lower} 
captures the sharp asymptotic rate.
\end{proposition}

\begin{proof}
The lower bound was established in Proposition~\ref{prop:poincare-gap}. 
We prove the upper bound.

Let $u \in W^{\alpha,2}_\Delta(\mathbb{T}_\delta)$ with 
$u_{\mathbb{T}_\delta} = 0$. Write $\mathbb{T}_\delta = I_1 \cup I_2$ with 
$I_1 = [0,1]$ and $I_2 = [1+\delta, 2+\delta]$, and let 
$u_i := u_{I_i}$ denote the average on $I_i$. Since 
$\mu_\Delta(I_1) = \mu_\Delta(I_2) = 1$ and $\mu_\Delta(\mathbb{T}_\delta) = 2$, 
the global vanishing average condition $u_{\mathbb{T}_\delta} = 0$ gives
\[
u_1 + u_2 = 0.
\]

\smallskip
\noindent\textbf{Step 1: $L^2$ decomposition.} Writing 
$u = (u - u_i) + u_i$ on each $I_i$,
\[
\|u\|^2_{L^2_\Delta(\mathbb{T}_\delta)} 
= \sum_{i=1}^2 \|u - u_i\|^2_{L^2(I_i)} + \sum_{i=1}^2 |u_i|^2 \, \mu_\Delta(I_i)
= \sum_{i=1}^2 \|u - u_i\|^2_{L^2(I_i)} + 2|u_1|^2,
\]
where the last equality uses $|u_1|^2 = |u_2|^2$.

\smallskip
\noindent\textbf{Step 2: control of intra-component oscillation.} On each 
interval $I_i$ of length~$1$ (independent of $\delta$), the classical 
fractional Poincar\'e inequality on bounded intervals gives
\[
\|u - u_i\|^2_{L^2(I_i)} 
\le C_1(\alpha) \, [u]^2_{W^{\alpha,2}(I_i)} 
\le C_1(\alpha) \, [u]^2_{W^{\alpha,2}_\Delta(\mathbb{T}_\delta)},
\]
where $C_1(\alpha)$ is independent of $\delta$.

\smallskip
\noindent\textbf{Step 3: control of the average.} Using $u_2 = -u_1$ and 
Jensen's inequality on $I_1 \times I_2$,
\[
|u_1|^2 = \frac{|u_1 - u_2|^2}{4} 
\le \frac{1}{4} \iint_{I_1 \times I_2} |u(t) - u(s)|^2 \, dt \, ds.
\]
For $(t,s) \in I_1 \times I_2$ and $\delta \ge 1$, one has 
$\delta \le |t-s| \le 2 + \delta \le 3\delta$, so
\[
|t-s|^{1+2\alpha} \le (3\delta)^{1+2\alpha} = 3^{1+2\alpha} \delta^{1+2\alpha},
\]
and therefore
\[
\iint_{I_1 \times I_2} |u(t) - u(s)|^2 \, dt \, ds 
\le 3^{1+2\alpha} \delta^{1+2\alpha} 
\iint_{I_1 \times I_2} \frac{|u(t) - u(s)|^2}{|t-s|^{1+2\alpha}} \, dt \, ds 
\le 3^{1+2\alpha} \delta^{1+2\alpha} \, [u]^2_{W^{\alpha,2}_\Delta(\mathbb{T}_\delta)}.
\]
Hence
\[
2|u_1|^2 \le \frac{3^{1+2\alpha}}{2} \delta^{1+2\alpha} \, 
[u]^2_{W^{\alpha,2}_\Delta(\mathbb{T}_\delta)}.
\]

\smallskip
\noindent\textbf{Step 4: combination.} Summing the bounds of Steps~2 and~3,
\[
\|u\|^2_{L^2_\Delta(\mathbb{T}_\delta)} 
\le \Big( 2 C_1(\alpha) + \frac{3^{1+2\alpha}}{2} \delta^{1+2\alpha} \Big) 
\, [u]^2_{W^{\alpha,2}_\Delta(\mathbb{T}_\delta)}.
\]
For $\delta \ge 1$, the right-hand side is bounded by 
$C^*(\alpha) \, \delta^{1+2\alpha} \, [u]^2_{W^{\alpha,2}_\Delta(\mathbb{T}_\delta)}$ 
with $C^*(\alpha) := 2 C_1(\alpha) + \tfrac{1}{2} 3^{1+2\alpha}$. Taking square 
roots and the supremum over admissible $u$ yields
\[
C_P(\delta) \le \big(C^*(\alpha)\big)^{1/2} \, \delta^{(1+2\alpha)/2},
\]
which proves the upper bound in \eqref{eq:two-sided-poincare}.
\end{proof}

\end{document}
